\section{Endpoint sufficiency for the finite mixture}
\label{app:mixture-endpoint}

\subsection{Pointwise and almost-sure formulations}
Fix $0<\pi_0<1$. Positivity gives a continuous posterior
$q(z)=\sigmoid(\logit\pi_0+\ell(z))$. Thus statistical sufficiency of
$S=\boldsymbol 1^{\mathsf T}Z$ is equivalent to
$\ell(z)=b(S)$ almost everywhere for some measurable $b$. The mixture density
$(1-\pi_0)m_0+\pi_0m_1$ is everywhere positive, so its null sets and Lebesgue
null sets agree.

Here the almost-everywhere qualifier can be removed. Use linear coordinates
$(s,u)\in\R\times\R^{T-1}$ with
$z=(u_1,\ldots,u_{T-1},s-\sum_{i<T}u_i)$, and write the transformed continuous
function as $F(s,u)$. Fubini's theorem implies that, for almost every $s$,
$F(s,u)=b(s)$ for almost every $u$. Continuity in $u$ extends this equality to
every $u$ for each such $s$. In particular,
$F(s,u)=F(s,0)$ on those fibers. Their $s$-values form a dense set, so continuity
in $s$ implies $F(s,u)=F(s,0)$ for every $(s,u)$. Hence
$\widetilde b(s)=F(s,0)$ supplies a $C^1$ pointwise version.

If $\ell(z)=\widetilde b(\sum_tz_t)$, the chain rule gives
\eqref{eq:mixture-gradient-test}. Conversely, that equation makes
$\nabla\ell(z)$ proportional to $\boldsymbol 1$. For any $z,z'$ on the same
sum hyperplane, the segment between them remains on that hyperplane and
\[
 \ell(z')-\ell(z)
 =\int_0^1\nabla\ell(z+u(z'-z))^{\mathsf T}(z'-z)\,du=0.
\]
All sum hyperplanes are connected. This proves
Proposition~\ref{prop:mixture-endpoint-gradient}, including the equivalence to
statistical sufficiency rather than merely a pointwise algebraic condition.
On a restricted support, connectedness of the relevant fibers must be
checked; the common-tail proposition additionally needs the stated full
unbounded support.

\subsection{Why the score improvement is strict}
Let $q_{\rm H}=\Prob(Y=1\mid Z)$ and
$q_{\rm E}=\Prob(Y=1\mid S)=\E(q_{\rm H}\mid S)$. Failure of the gradient
condition implies, by the preceding equivalence, that $q_{\rm H}$ is not
$S$-measurable up to a null set. Therefore
\[
 \E[(q_{\rm H}-q_{\rm E})^2]>0.
\]
Conditional expectation gives the Brier-risk difference exactly as this
quantity. Likewise, the logarithmic-risk difference equals
\[
 \E\!\left[
 D_{\rm KL}\{\operatorname{Bern}(q_{\rm H})
           \Vert\operatorname{Bern}(q_{\rm E})\}\right]
 = I(Y;Z\mid S)>0.
\]
Both optimal logarithmic risks are conditional entropies and are finite,
bounded by $H(Y)\leq\log2$. These are strict population comparisons between
optimal forecasts under one specified model. They are not guarantees for a
fitted classifier under an unknown data-generating distribution.

One can also see why a same-endpoint counterexample is not only a null-set
observation. Two unequal continuous values on one fiber persist in disjoint
neighborhoods and in nearby fibers. Positive densities assign positive
conditional mass to those fiber neighborhoods for almost every nearby sum.
The posterior then has positive conditional variance on a set of sums with
positive probability. The Fubini argument above formalizes this extension.

\subsection{The two levels of mixture responsibility}
All sums are finite, so differentiation commutes with summation without a
domination condition. Since the gates depend on the fixed design and not on
$z_t$,
\[
 \partial_{z_t}\log f_{yjt}(z_t)
 =\sum_k\frac{\rho_{jkt}g_{yjkt}(z_t)}{f_{yjt}(z_t)}
          \partial_{z_t}\log g_{yjkt}(z_t).
\]
Write $\tau_{yj}(z)=\varpi_j\prod_t f_{yjt}(z_t)/m_y(z)$ and
$\eta_{yjkt}(z_t)=\rho_{jkt}g_{yjkt}(z_t)/f_{yjt}(z_t)$.
Differentiate the product within each shared-state likelihood and then the
outer state mixture to obtain
\begin{equation}
 \partial_{z_t}\log m_y(z)
 =\sum_j\tau_{yj}(z)\sum_k\eta_{yjkt}(z_t)
       \partial_{z_t}\log g_{yjkt}(z_t).
 \label{eq:mixture-score-responsibilities}
\end{equation}
This is the elementary conditional-score identity of incomplete-data
likelihood calculations \cite{louis1982observed}, differentiated with respect
to an observation coordinate. The Gaussian component score is
$-(z_t-\mu)/\sigma^2$. For a Student component with location $\mu$, scale $b$,
and degrees of freedom $\nu$, the score is
\[
 -\frac{(\nu+1)(z_t-\mu)}{\nu b^2+(z_t-\mu)^2}.
\]
The checker evaluates both responsibility levels explicitly and compares the
resulting derivative with central differences of the production marginal
likelihood at fixed, unfitted parameter values.

\subsection{Independent-increment specialization}
\begin{proposition}[Common exponential tilt]
\label{prop:mixture-independent-tilt}
Let $T\geq2$ and suppose
$m_y(z)=\prod_t f_{yt}(z_t)$, where every $f_{yt}$ is a strictly positive $C^1$
density on $\R$. Then $\sum_tZ_t$ is sufficient for $Y$ if and only if there
exists one constant $c$ such that, for every $t$,
\begin{equation}
 f_{1t}(x)=\exp\{cx-\psi_t(c)\}f_{0t}(x),\qquad
 \psi_t(c)=\log\int_{\R}e^{cu}f_{0t}(u)\,du<\infty.
 \label{eq:mixture-common-tilt}
\end{equation}
The slope $c$ is common to all increments; the normalizing intercepts may
differ.
\end{proposition}
\begin{proof}
Write $h_t(x)=\log\{f_{1t}(x)/f_{0t}(x)\}$. The log Bayes factor is
$\sum_t h_t(z_t)$. By Proposition~\ref{prop:mixture-endpoint-gradient},
$h'_r(x)=h'_t(u)$ for all $x,u$ and any distinct $r,t$. Because $T\geq2$ and
the support is all of $\R^T$, one may vary either coordinate while holding the
other fixed. Every derivative must therefore be the same constant $c$.
Integration gives $h_t(x)=cx+b_t$; normalization of $f_{1t}$ implies
$b_t=-\psi_t(c)$ and finiteness of the integral. Conversely,
\eqref{eq:mixture-common-tilt} gives
$\ell(z)=c\sum_tz_t-\sum_t\psi_t(c)$, a function of the endpoint.
\end{proof}

If every one-step likelihood ratio tends to one at both tails, then
$cx+b_t\to0$ forces $c=0$ and $b_t=0$. Consequently endpoint sufficiency in
that independent-increment family requires identical outcome densities at
every step. This argument concerns a fixed sample dimension and a binary
experiment; it does not invoke an unrestricted fixed-dimensional-sufficiency
characterization for an entire parametric family.

\subsection{Proof of the common-tail proposition}
For any fixed $s\in\R$ and $R>0$, let
\[
 z_R=(R,\ldots,R,s-(T-1)R).
\]
Every coordinate tends to one of the two tails as $R\to\infty$, whereas its
sum remains $s$. For each catalogue state,
\[
 A_j(R)=\prod_{t=1}^T\frac{f_{1jt}(z_{R,t})}{f_{0jt}(z_{R,t})}\longrightarrow1.
\]
Writing the marginal likelihood ratio as a convex combination yields
\[
 \frac{m_1(z_R)}{m_0(z_R)}
 =\sum_j\frac{\varpi_j\prod_t f_{0jt}(z_{R,t})}{m_0(z_R)}A_j(R),\qquad
 \min_j A_j(R)\leq\frac{m_1(z_R)}{m_0(z_R)}\leq\max_j A_j(R).
\]
The finite catalogue makes both bounds tend to one. If $S$ were sufficient,
the pointwise version proved above would make this ratio equal to
$\exp\{b(s)\}$ for every $R$, hence $b(s)=0$. Since $s$ was arbitrary,
$m_1=m_0$. The converse is immediate. The strict-score statement follows from
the earlier risk identities.

The proof requires common outcome-specific catalogue masses. It does not
justify an exchange of a tail limit and an arbitrary continuous nuisance
integral. Nor does it apply unchanged after censoring, truncation, or clipping
has replaced the full-support observation law.

For a Student component with a fixed finite location, positive scale $b$, and
degrees of freedom $\nu$, its density has leading term
$c_\nu b^\nu |x|^{-(\nu+1)}$ at either tail. If its mixing weight, scale, and
$\nu$ do not depend on the outcome, that term is identical for both outcomes;
finite Gaussian components decay faster. This proves
\eqref{eq:mixture-common-tail} for the stated empirical and synthetic
catalogues, including the shifted Student component in the latter.

Finally, the empirical mean calculation does not assume a favorable fitted
weight vector. Its Student component is centered, so at normalized volume
$r>0$ the sign of the directional mean for $a>0$ is the sign of
\[
 w(1+r)^{1/2}-\frac{1-w}{4}(1+r)^{3/10}
 =(1+r)^{3/10}\left[w(1+r)^{1/5}-\frac{1-w}{4}\right].
\]
For $w\geq1/5$, the bracket is strictly positive when $r>0$. The catalogue
contains positive amplitudes and its smoothed masses are positive, giving
opposite nonzero marginal means under the two outcomes at any such time.
The outcome laws are therefore distinct. If all volumes are zero, all
location shifts vanish and this informativeness argument does not apply.
